\section{Problem Formulation}

We formalize the problem of post-admission memory validity revision. Let a scenario be defined as a tuple:
\begin{equation}
x = (M, E, q, c),
\end{equation}
where $M=\{m_i\}_{i=1}^{n}$ is the persistent memory store, $E=\{e_j\}_{j=1}^{k}$ is the chronological evidence ledger, $q$ is the downstream query, and $c$ contains public policy and scope constraints. Both memory records and evidence events are associated with unique, stable identifiers and natural-language payloads.

The public revision view is constructed via a deterministic function $B$:
\begin{equation}
V = B(M, E, q, c, s^0, \Gamma),
\end{equation}
where $s^0$ is the initial state map of the memory records, and $\Gamma$ is the published transition policy. Crucially, $V$ excludes gold labels, private metadata, and evaluator-specific information, exposing only the candidate records, evidence ledger, and active policies to the agent.

The memory state space $\Sigma$ is defined over five distinct epistemic classes:
\begin{equation}
\begin{split}
\Sigma=\{&\mathsf{Current},\mathsf{Blocked},\mathsf{Unresolved},\\
&\mathsf{OutOfScope},\mathsf{ShouldNotStore}\}.
\end{split}
\end{equation}
The downstream reuse license is a function $\lambda: \Sigma \rightarrow \{0, 1\}$:
\begin{equation}
\lambda(z) = \mathbf{1}[z = \mathsf{Current}],
\end{equation}
which dictates that a memory record $m_i$ is licensed for future reuse in downstream tasks if and only if its validity state $s_i$ is $\mathsf{Current}$. Records with other states (e.g., blocked due to unmet dependencies, unresolved due to conflicting evidence, contextually out-of-scope, or excluded by safety policy) are restricted from reuse.

\paragraph{Output Contract.}
The agent's output is structured as a tuple:
\begin{equation}
r = (d, \mathbf{s}, Z, f, a),
\end{equation}
where $d$ is a downstream decision (e.g., use memory, escalate, ask clarification, refuse due to policy), $\mathbf{s}: M \rightarrow \Sigma$ is the final memory validity mapping, $Z \subseteq I_E$ is the response-level evidence set, $f \in \mathcal{F}$ is the failure diagnosis, and $a$ is the final generated answer. For a committed patch, provenance is finer grained: $\widehat Z:M\rightarrow 2^{I_E}$ associates each changed record with its own evidence witness.

\paragraph{Diagnostic Measurement Model.}
To assess agent performance without hiding failure modes inside a single conjunction, we define seven response channels: decision, state, evidence, minimality, diagnosis, answer, and schema. They correspond to decision accuracy, memory-state accuracy, evidence $F_1$, minimal-evidence exact match, diagnosis accuracy, answer key-fact accuracy, and schema compliance. We probe these channels with seven controlled corruptions $\mathcal{O}$: wrong decision, wrong state, wrong evidence, overcitation, wrong diagnosis, malformed schema, and missing answer. Trace audit is evaluated separately because replayability is a property of an executed transition, not of a five-field response.

For each scenario, we define an error vector $\mathbf{e}(r, g) = (\ell_k(r, g))_{k \in \mathcal{K}}$, where each $\ell_k$ measures the correctness of output channel $k$ against gold label $g$. The diagnostic profile of an agent $h$ is the expected error vector:
\begin{equation}
\mathbf{R}(h) = \mathbb{E}_{\xi} [\mathbf{e}(h(B(\xi)), g(\xi))].
\end{equation}
The sensitivity matrix $A=(a_{kj})\in\mathbb{R}^{7\times7}$ records the decrement in metric $k$ under corruption $o_j$. All operators, mappings, and calibration constants are fixed exclusively on the L3 development split before L4 evaluation. Full rank is an empirical audit of this chosen corruption basis, not a claim that all real model failures are linearly separable. We report the complete error vector as the primary outcome. Joint revision success, $J(\mathbf b)=\prod_k(1-b_k)$, is retained only as a secondary complete-case statistic because it maps many distinct error signatures to the same zero value.
